\subsection{系数在特征为 p 的环中的 Witt 向量环}

命题 5．——设 A 为特征为 p 的环（A，V，第 2 页）。对 W(A) 中的任意元素 a、b 以及正整数 m、n，若 \(\boldsymbol{a} = (a_{n})_{n \in \mathbb{N}}\)，

\begin{enumerate}
\def\labelenumi{(\arabic{enumi})}
\setcounter{enumi}{50}
\tightlist
\item
\end{enumerate}

\[
\mathrm{F} (\boldsymbol {a}) = (a _ {n} ^ {p}) _ {n \in \mathbf {N}}\tag{52}
\]

\[
p. \boldsymbol {a} = \operatorname{VF} (\boldsymbol {a}) = \operatorname{FV} (\boldsymbol {a}) = (0, a _ {0} ^ {p}, a _ {1} ^ {p}, \dots)\tag{53}
\]

\[
\mathrm{V} ^ {m} (a) \times \mathrm{V} ^ {n} (b) = \mathrm{V} ^ {m + n} (\mathrm{F} ^ {n} (a) \times \mathrm{F} ^ {m} (b)).
\]

公式 (51) 由第 3 节例 4 得到。由此立即推出等式

\[
\operatorname{VF} (\boldsymbol {a}) = \operatorname{FV} (\boldsymbol {a}) = (0, a _ {0} ^ {p}, a _ {1} ^ {p}, \dots),
\]

且等式 \(p\) . \(a = \text{FV}(a)\) 已由第 5 节命题 3 证明，从而得到 (52)。

证明 (53)。由公式 (37)（其中以 \(V^n(\boldsymbol{b})\) 代替 \(\boldsymbol{b}\)），有

\[
\mathrm{V} ^ {m} (\boldsymbol {a}) \times \mathrm{V} ^ {n} (\boldsymbol {b}) = \mathrm{V} ^ {m} \left(\boldsymbol {a} \times \mathrm{F} ^ {m} \left(\mathrm{V} ^ {n} (\boldsymbol {b})\right)\right).\tag{54}
\]

由公式 (37) 还可推出

\[
\mathrm{V} ^ {n} \left(\mathrm{F} ^ {m} (\boldsymbol {b})\right) \times \boldsymbol {a} = \mathrm{V} ^ {n} \left(\mathrm{F} ^ {m} (\boldsymbol {b}) \times \mathrm{F} ^ {n} (\boldsymbol {a})\right).\tag{55}
\]

于是，公式 (53) 由 (54)、(55) 以及关系 \(F^{m} \circ V^{n} = V^{n} \circ F^{m}\) 得到，而后一个关系本身是 (51) 的结果。

推论．——若 m、n 为两个正整数，则有

\[
\mathrm{V} _ {m} (\mathrm{A}) \times \mathrm{V} _ {n} (\mathrm{A}) \subset \mathrm{V} _ {m + n} (\mathrm{A}).
\]

这由公式 (53) 得到，因为 \(V_{m}(A)\) 是 \(V^{m}:W(A)\to W(A)\) 的像。

命题 6．——设 A 为一个环。

\begin{enumerate}
\def\labelenumi{\alph{enumi})}
\item
  对每个整数 \(k \geqslant 1\)，都有 \((\mathrm{V}_1(\mathrm{A}))^k = p^{k-1} \cdot \mathrm{V}_1(\mathrm{A})\)。
\item
  假设 A 是特征为 p 的环。在环 W(A) 上，V₁(A)-进拓扑与 p-进拓扑重合，并且二者都比乘积拓扑 ℰ 精细（参见第 6 节）。环 W(A) 关于 p-进拓扑是分离且完备的。
\end{enumerate}

证明 \(a\)：对 \(k\) 归纳。当 \(k = 1\) 时显然成立。设 \(k \geqslant 2\)。根据归纳假设，有 \(\mathrm{V}_{1}(\mathrm{A})^{k-1} = p^{k-2} \cdot \mathrm{V}_{1}(\mathrm{A})\)，从而 \(\mathrm{V}_{1}(\mathrm{A})^{k} = p^{k-2} \cdot (\mathrm{V}_{1}(\mathrm{A}))^{2}\)。但由第 5 节命题 3 的 \(d\)、公式 (31)，有 \((\mathrm{V}_{1}(\mathrm{A}))^{2} = p \cdot \mathrm{V}_{1}(\mathrm{A})\)，故得 \(a\)。

现在假设 A 具有特征 p。由于有

\[
p. \mathrm{W} (\mathrm{A}) = \mathrm{VF} (\mathrm{W} (\mathrm{A})) \subset \mathrm{V} _ {1} (\mathrm{A}) \quad (\text { formula (52) }),
\]

由 \(a\) 得到包含关系 \(p^{k} \cdot \mathbf{W}(\mathbf{A}) \subset (\mathbf{V}_{1}(\mathbf{A}))^{k} \subset p^{k-1} \cdot \mathbf{W}(\mathbf{A})\)，并由命题 5 的推论得到包含关系 \((\mathbf{V}_{1}(\mathbf{A}))^{k} \subset \mathbf{V}_{k}(\mathbf{A})\)，对每个整数 \(k \geqslant 1\) 均如此。由此得到 \(b\) 的第一条断言。

设 \(k\) 为一个 \(\geqslant 1\) 的整数。根据公式 (52)，W(A) 的理想 \(p^{k}\) . W(A) 是 W(A) 中元素 \(\boldsymbol{a} = (a_{n})_{n \in \mathbf{N}}\) 所成的集合，其中元素满足 \(a_{n} = 0\)（\(n < k\)）且 \(a_{n} \in A^{p^{k}}\)（\(n \geqslant k\)）。因此它关于拓扑 \(\mathcal{G}\) 是闭的。由于 W(A) 关于拓扑 \(\mathcal{G}\) 分离且完备，而 W(A) 的理想 \(p^{k}\) . W(A)（对 \(k \geqslant 1\)）构成 W(A) 中 \(p\)-进拓扑的 0 的邻域基，故环 W(A) 关于 \(p\)-进拓扑分离且完备（TG，III，第 26 页，命题 10 的推论 1）。

命题 7．——设 A 为特征为 p 的完美环。

\begin{enumerate}
\def\labelenumi{\alph{enumi})}
\item
  对 W(A) 的每个元素 \(\boldsymbol{a} = (a_{n})_{n \in \mathbb{N}}\)，通项为 \(p^{n}\tau(a_{n}^{p^{-n}})\) 的级数在 W(A) 中收敛，且其和为 a。
\item
  在 \(\mathbf{W}(\mathbf{A})\) 上，\(\mathbf{V}_1(\mathbf{A})\)-进拓扑、\(p\)-进拓扑和拓扑 \(\mathcal{C}\) 重合。更准确地说，有 \(\mathbf{V}_n(\mathbf{A}) = p^n\) . \(\mathbf{W}(\mathbf{A}) = (\mathbf{V}_1(\mathbf{A}))^n\)，对每个整数 \(n \geqslant 0\) 均成立。特别地，\(\Phi_0\) 定义了从 \(\mathbf{W}(\mathbf{A})/p\) . \(\mathbf{W}(\mathbf{A})\) 到 \(\mathbf{A}\) 的同构。
\end{enumerate}

根据定义（A，V，第 5 页），映射 \(a \mapsto a^{p}\) 是环 A 的自同构。因此由命题 5，F 是环 W(A) 的自同构，并且对每个 \(n \in \mathbf{N}\)，有

\[
p ^ {n}. \mathrm{W} (\mathrm{A}) = \mathrm{V} ^ {n} \mathrm{F} ^ {n} (\mathrm{W} (\mathrm{A})) = \mathrm{V} ^ {n} (\mathrm{W} (\mathrm{A})) = \mathrm{V} _ {n} (\mathrm{A}).
\]

特别地，有 \((\mathbf{V}_1(\mathbf{A}))^n = (p.\mathbf{W}(\mathbf{A}))^n = p^n.\mathbf{W}(\mathbf{A})\)。断言 \(b)\) 由此得到。由命题 5，有

\[
p ^ {n} \cdot \tau (a _ {n} ^ {p ^ {- n}}) = \mathrm{V} ^ {n} \mathrm{F} ^ {n} \tau (a _ {n} ^ {p ^ {- n}}) = \mathrm{V} ^ {n} \tau (a _ {n}),
\]

从而由第 6 节命题 4 得到断言 a)。

命题 8．——设 A 为特征为 p 的域。环 W(A) 是一个局部整环，分离且完备，其极大理想为 V₁(A)，剩余域同构于 A。若域 A 是完美的，则环 W(A) 是一个离散赋值环，其极大理想为 p.W(A)。

同态 \(\Phi_{0}\) 定义了从 W(A)/V \(_{1}\) (A) 到 A 的同构（第 7 节例 1）。因此，V \(_{1}\) (A) 是 W(A) 的极大理想。由于环 W(A) 关于 V \(_{1}\) (A)-进拓扑分离且完备，它是一个局部环，极大理想为 V \(_{1}\) (A)（命题 6 的 b)，以及 III，第 2 节第 13 号命题 19）。

设 \(\pmb{a}\) 和 \(\pmb{b}\) 是 \(\mathbf{W}(\mathbf{A})\) 的两个非零元素。存在整数 \(m \geqslant 0\) 和 \(n \geqslant 0\)，以及元素 \(\pmb{a}' = (a_n')_{n \in \mathbb{N}}\) 和 \(\pmb{b}' = (b_n')_{n \in \mathbb{N}}\)，它们属于 \(\mathbf{W}(\mathbf{A})\)，使得 \(\pmb{a} = \mathbf{V}^m(\pmb{a}')\)、\(\pmb{b} = \mathbf{V}^n(\pmb{b}')\)，且 \(a_0'\)、\(b_0'\) 是 \(\mathbf{A}\) 中的非零元素。于是，下标为 \(m + n\) 的 \(\pmb{a} \times \pmb{b}\) 的分量等于下标为 0 的 \(\mathbf{F}^n(\pmb{a}') \times \mathbf{F}^m(\pmb{b}')\) 的分量（公式 (53)），即等于 \(a_0'^{p^n} b_0'^{p^m}\)（公式 (51) 及第 3 节例 2）。因此 \(\pmb{a} \times \pmb{b}\) 非零，且 \(\mathbf{W}(\mathbf{A})\) 是整环。

若域 A 是完美的，则极大理想 \(V_{1}(A)\) 属于 \(W(A)\)，并等于 p.\(W(A)\)

（命题 7 的 b)），从而 W(A) 是一个离散赋值环（VI，第 3 节第 6 号命题 9 的 c)）。

备注．——1) 设 A 为特征为 p 的域。可以证明，环 W(A) 是诺特环当且仅当 A 是完美的（第 43 页习题 9）。

\begin{enumerate}
\def\labelenumi{\arabic{enumi})}
\setcounter{enumi}{1}
\tightlist
\item
  设 A 为特征为 p 的环。由命题 5，有公式
\end{enumerate}

\[
\mathrm{F} _ {m} ^ {n} [ a _ {0}, \dots , a _ {n + m - 1} ] = [ a _ {0} ^ {p ^ {n}}, \dots , a _ {m - 1} ^ {p ^ {n}} ]
\]

\[
p ^ {n} \cdot [ a _ {0}, \dots , a _ {n + m - 1} ] = [ \underbrace {0 , \dots , 0} _ {n \text {   times }}, a _ {0} ^ {p ^ {n}}, \dots , a _ {m - 1} ^ {p ^ {n}} ]
\]

对每个 Witt 向量 \([a_0, \dots, a_{n+m-1}]\)，其长度为 \(n + m\)，均成立。

事实上，映射 \(F:W(A)\to W(A)\) 使我们能够通过对 \(V_{m}(A)\) 取商来定义映射 \(\overline{F}_{m}:W_{m}(A)\to W_{m}(A)\)。有公式

\[
\overline {{{{\mathrm{F}}}}} _ {m} \left[ a _ {0}, \dots , a _ {m - 1} \right] = \left[ a _ {0} ^ {p}, \dots , a _ {m - 1} ^ {p} \right].
\]

映射 \(V_{m}^{1} \circ \overline{F}_{m}\) 和 \(\overline{F}_{m+1} \circ V_{m}^{1}\) 从 \(W_{m}(A)\) 到 \(W_{m+1}(A)\)，彼此相等，并且通过取商由 \(W_{m+1}(A)\) 中乘以 p 得到。

